\section{Problem Definition and Game-Theoretic Formulation}
\label{sec_model}

We model the dispute as a simultaneous $2\times2$ non-cooperative game between two aggregated actors. The first actor is the offshore coalition, which includes developers and institutional supporters of the project. The second actor is the surf-tourism coalition, which includes surfers, tourism operators, local businesses, and public or civic actors whose income depends on preserving coastal and wave conditions. This abstraction follows the logic of normal-form games used to study strategic interaction under conflict and incomplete alignment of interests \citep{gillman2019,sobhee_game_2017}.

The offshore coalition chooses between \textit{Implement} and \textit{Withdraw}. The surf-tourism coalition chooses between \textit{Accept} and \textit{Fight}. The latter action stands for organized resistance through public hearings, legal dispute, technical contestation, and negotiation pressure. The model is not meant to deny the diversity within each group. It is meant to capture the strategic bottleneck that appears once both sides must decide whether they escalate or accommodate.

\subsection{Variables and economic meaning}

The worksheet defines the following variables:

\begin{itemize}
\item $R$: expected gross return from the offshore project.
\item $LR$: expected reduction in offshore return due to concessions, redesign, delay, or loss associated with the dispute.
\item $C_s$: offshore-side cost of sustaining the dispute.
\item $I$: current income of the surf-tourism side in the status quo.
\item $L$: expected loss to surf-tourism income if the project is implemented.
\item $C_a$: cost incurred by the surf-tourism side when it fights.
\item $G$: compensation expected by surf-tourism actors if the project proceeds.
\item $P$: probability that the offshore side prevails in the dispute.
\item $\alpha$: residual fraction of contestation cost borne by the surf-tourism side if the offshore coalition withdraws.
\end{itemize}

The first four payoffs are direct. If the project is implemented and accepted, the offshore side obtains $R$ and the surf-tourism side retains $I-L$. If the project is withdrawn and accepted, the offshore side obtains $0$ and the surf-tourism side retains $I$. If the offshore side withdraws after opposition has already been mobilized, the worksheet assumes a residual contestation cost $\alpha C_a$ for the surf-tourism side, giving the payoff $I-\alpha C_a$.

\subsection{Expected payoff under conflict}

Conflict arises only in the cell where the offshore coalition implements and the surf-tourism coalition fights. In that cell, payoffs depend on the probability of prevailing:

\begin{equation}
G_1 = P(R-LR-C_s-G) + (1-P)(-C_s),
\label{eq:g1}
\end{equation}

\begin{equation}
G_2 = (1-P)(I-C_a) + P(I-L+G-C_a).
\label{eq:g2}
\end{equation}

Equation \ref{eq:g1} states that if offshore proponents win, they keep the project return net of expected concession/loss $LR$, legal or political cost $C_s$, and compensation $G$; if they lose, they only absorb the dispute cost. Equation \ref{eq:g2} states that if surf-tourism actors win, they preserve the status-quo income net of contestation cost; if they lose, they bear the local loss $L$ but recover part of it through compensation $G$.

The resulting payoff matrix is shown in Table \ref{tab:abstract_game}.

\begin{table}[H]
\centering
\caption{Abstract decision problem for offshore wind versus surf-tourism.}
\label{tab:abstract_game}
\begin{tabular}{lcc}
\toprule
 & Accept & Fight \\
\midrule
Implement & $(R, I-L)$ & $(G_1, G_2)$ \\
Withdraw & $(0, I)$ & $(0, I-\alpha C_a)$ \\
\bottomrule
\end{tabular}
\end{table}

\subsection{Dominance conditions}

The offshore coalition prefers implementation to withdrawal whenever the conflict payoff remains positive:

\begin{equation}
G_1 > 0 \Longleftrightarrow P(R-LR-G) > C_s.
\label{eq:offshore_condition}
\end{equation}

Under maximum uncertainty, $P=0.5$, this becomes

\begin{equation}
R-LR-G > 2C_s.
\label{eq:offshore_condition_half}
\end{equation}

This condition is substantively important because it shows that the offshore side does not compare $R$ with zero in a naive way. It compares a probability-weighted net return with the direct cost of sustaining conflict.

For the surf-tourism coalition, fighting is preferred to accepting when the expected conflict payoff is at least as high as passive acceptance:

\begin{equation}
G_2 \geq I-L \Longleftrightarrow C_a \leq (1-P)L + PG.
\label{eq:opponent_condition}
\end{equation}

With $P=0.5$, the threshold simplifies to

\begin{equation}
2C_a \leq L + G.
\label{eq:opponent_condition_half}
\end{equation}

This is the key behavioral result from the worksheet. Local actors are more likely to resist when expected losses are salient, the cost of mobilization is manageable, and compensation is sufficiently large to soften the downside risk of losing. The last term may seem counterintuitive, but it is strategically coherent: a compensation scheme can simultaneously make implementation politically easier and resistance less dangerous for the local side.

\subsection{Baseline calibration from the worksheet}

The worksheet proposes the baseline values
\[
R=100,\quad LR=20,\quad C_s=3,\quad I=10,\quad L=3,\quad \alpha=0,\quad C_a=3,\quad G=4,\quad P=0.5.
\]

Substituting these values into Equations \ref{eq:g1} and \ref{eq:g2} yields
\[
G_1 = 35
\]
and
\[
G_2 = 7.5.
\]

Table \ref{tab:baseline_model} gives the calibrated matrix.

\begin{table}[H]
\centering
\caption{Calibrated game using the worksheet baseline.}
\label{tab:baseline_model}
\begin{tabular}{lcc}
\toprule
 & Accept & Fight \\
\midrule
Implement & $(100, 7)$ & $(35, 7.5)$ \\
Withdraw & $(0, 10)$ & $(0, 10)$ \\
\bottomrule
\end{tabular}
\end{table}

The offshore side has a dominant preference for implementation because both $100>0$ and $35>0$. The surf-tourism side prefers fighting when implementation is expected because $7.5>7$. With $\alpha=0$, accepting and fighting are equivalent when the project is withdrawn because mobilization costs do not survive the withdrawal decision. The resulting pure-strategy equilibrium is therefore \textit{Implement/Fight}.

\subsection{Nash equilibrium, Pareto ranking, and negotiation meaning}

The calibrated game reproduces a classic policy problem. The equilibrium cell is not the collectively best cell. Summing the two players' payoffs, \textit{Implement/Accept} yields $107$, while \textit{Implement/Fight} yields only $42.5$. The Nash outcome is therefore conflictual and Pareto-inferior to negotiated acceptance. This divergence is analytically useful because it captures the political reality observed in maritime planning: mutually recognized value can coexist with persistent conflict when bargaining institutions are weak or compensation rules are incomplete.

This also helps connect the present model to the support review of game-theory applications in coastal governance. In tragedy-of-the-commons settings, sanctioning and subsidizing instruments are used to move the equilibrium away from individually rational but collectively damaging behavior \citep{hardin1968tragedy,ostrom_2015,sobhee_game_2017}. In the present litigation-oriented game, the analogue is to redesign $LR$, $G$, $C_s$, and $C_a$ through policy. Better siting and technical redesign reduce $LR$ and $L$. Early mediation and procedural clarity can reduce $C_s$ and $C_a$. Compensation redesign changes $G$, but it must be handled carefully because compensation that only insures the downside of resistance may preserve the conflict equilibrium instead of dissolving it.

\subsection{Methodological implications}

The model is intentionally parsimonious. It omits repeated interaction, intra-coalition heterogeneity, and dynamic reputation effects, all of which matter in real maritime negotiations. Even so, it is already useful for three reasons. First, it makes the decision problem explicit enough to be discussed by non-specialists in policy and local governance. Second, it shows exactly which quantities have to be estimated or negotiated for the conflict to change shape. Third, it provides a reproducible bridge between the qualitative narrative of Ericeira and the stylized calibration based on Peniche, which is the role later operationalized in the R workflow delivered with this paper.
